\section{Data}

The data consist of a panel of 19 euro area member states observed at monthly frequency from
January 1998 to December 2025, assembled from Eurostat and the ECB Data Portal.
Data for Croatia are available from January 1999; all remaining members are observed over
the full period, yielding 6,384 country-month observations in total.
Table~\ref{tab:summary} presents summary statistics.
Table~\ref{tab:sources} documents data sources.
Table~\ref{tab:coverage} lists country coverage with sample start and end dates and
observation counts.

The central data contribution of this paper is the construction of country-specific tradable and
non-tradable inflation measures directly from disaggregated HICP components following the
European Classification of Individual Consumption according to Purpose (ECOICOP).
For each country $i$ and ECOICOP division $j$, price indices are downloaded from Eurostat
dataset \texttt{prc\_hicp\_midx} (monthly, index base 2015\,=\,100) and annual expenditure
shares from \texttt{prc\_hicp\_inw}.
The year-over-year inflation rate for component $j$ is
$\pi_{ijt} = 100\times(P_{ijt}/P_{ij,t-12}-1)$,
where $P_{ijt}$ denotes the price level index.
Tradable and non-tradable inflation are then expenditure-weighted averages of component
inflation rates within each basket:
\begin{equation}
    \pi^k_{it} = \frac{\sum_{j \in \mathcal{k}} w_{ijt} \, \pi_{ijt}}{\sum_{j \in \mathcal{k}} w_{ijt}},
    \quad k \in \{T, N\}
    \label{eq:basket}
\end{equation}
where $w_{ijt}$ is the expenditure share of component $j$ in country $i$ at time $t$, expanded
to monthly frequency by assigning each month the annual weight of its calendar year.
The tradable basket $\mathcal{T}$ comprises CP01 (food and non-alcoholic beverages), CP02
(alcoholic beverages and tobacco), CP03 (clothing and footwear), CP05 (furnishings and
household equipment), and CP07 (transport).
The non-tradable basket $\mathcal{N}$ comprises CP04 (housing, water, electricity, and gas),
CP06 (health), CP08 (communication), CP09 (recreation and culture), CP10 (education),
CP11 (restaurants and hotels), and CP12 (miscellaneous goods and services).
Expenditure shares are normalized to sum to one across the twelve ECOICOP divisions.
Most existing work proxies tradable inflation using producer price indices or import prices
rather than constructing it directly from disaggregated consumer price components with
matching time-varying expenditure weights \citep{ha2023}; the present approach follows the
HICP-based decomposition of \citet{andriantomanga2025}, who applies the same
classification scheme to a broader sample of 40 countries covering the period from January
2000 to August 2022.

The analysis draws on four additional groups of variables.
The oil price shock series are from \citet{baumeister2019}: the market-based oil price
expectation shock captures forward-looking shifts in the expected path of oil prices, while the
structural oil supply shock isolates exogenous disruptions to global crude output identified
via sign and magnitude restrictions.
Global supply chain conditions enter through the Global Supply Chain Pressure Index (GSCPI)
of \citet{benigno2022}, which aggregates shipping costs, airfreight rates, and purchasing
managers' index backlogs into a single composite barometer of international logistics
disruptions.
Domestic cyclical conditions are captured by the seasonally adjusted unemployment rate
(percent of the active population) from Eurostat dataset \texttt{une\_rt\_m} and by the
year-over-year growth rate of the seasonally and calendar-adjusted manufacturing
production index from Eurostat dataset \texttt{sts\_inpr\_m}.
External monetary conditions enter through the log monthly change in the ECB nominal
effective exchange rate against 40 trading partners, retrieved from the ECB Data Portal API.
Import exposure is measured as imports of goods and services as a share of GDP from
Eurostat national accounts dataset \texttt{nama\_10\_gdp}, expanded to monthly frequency
by carrying forward each annual observation for the corresponding calendar year.
Summary statistics, data sources, and country coverage are reported in
Tables~\ref{tab:summary}--\ref{tab:coverage}.
